\section{October 2026: the newest questions, derived step by step}

\begin{summary}{What was asked in October 2026}
Two committees: Hoedt, Schlüter, Zoitl and Mayr, Schlüter, Luukkonen, and one with Zoitl, Hochreiter, Kobler.
\textsf{Schlüter:} what is deep learning, what is a layer, linear activations, sigmoid versus step function,
learning rate and batch size, SGD versus Adam, receptive field and pooling. \
\textsf{Hochreiter:} delta errors, vanishing gradients and their fixes, initialization and its theory, the spectral
norm and the structure of $W$, regularization, optimizers, activations, normalizations (RMSNorm), AdamW versus
Adam with L2, the SVM primal and the meaning of the ``1''. \
\textsf{Rule for every answer:} one sentence to learn by heart, then derive the formula on the board.
\end{summary}

\begin{qabox}[Schlüter]{\asked What is deep learning? What is a layer, and what happens in a layer?}
\from{DL Basic Techniques (DLNN I) · Ch.~2 Neural networks · asked October 2026}
\ans Deep learning trains a network of many layers end to end, and every layer turns the representation of
the layer below into a more useful one. Inside a layer there are two steps: a linear map, then a non-linear
activation.
\formula One neuron sums its weighted inputs, adds a bias and bends the result. Stacking the weight rows of all
neurons of a layer gives a matrix:
\[ s_i=\sum_j w_{ij}a_j+b_i \;\Longrightarrow\; \mathbf{s}=W\mathbf{a}+\mathbf{b},\qquad \mathbf{a}^{\text{new}}=f(\mathbf{s}). \]
\pic{A layer is ``linear, then bend''. Deep means features built from features: edges, then parts, then objects.}
\end{qabox}

\begin{qabox}[Schlüter]{\asked Can an activation function be linear?}
\from{DL Basic Techniques (DLNN I) · Ch.~2 · asked October 2026}
\ans At the output, yes, for regression. In the hidden layers, no, because linear times linear stays linear.
\formula Two layers without activation collapse into one:
\[ W_2(W_1\mathbf{x}+\mathbf{b}_1)+\mathbf{b}_2=\underbrace{W_2W_1}_{W'}\mathbf{x}+\underbrace{W_2\mathbf{b}_1+\mathbf{b}_2}_{\mathbf{b}'}. \]
So any depth gives one linear model, which cannot solve XOR. For regression the target is any real number, so
the output stays linear.
\end{qabox}

\begin{qabox}[Defense]{\asked Supervised or unsupervised: what is the difference, and why choose one?}
\from{ML: Supervised Techniques (Lehner) unit~1 · asked October 2026}
\ans Supervised learning has a label for every input and learns the map from input to label, unsupervised
learning has no labels and looks for structure in the inputs.
\formula Supervised: data $\{(\mathbf{x}_n,y_n)\}$, learn $g(\mathbf{x})\approx y$ (classification, regression).
Unsupervised: data $\{\mathbf{x}_n\}$, learn clusters, a density $p(\mathbf{x})$ or a projection (PCA).
Self-supervised: labels made from the data, for example the next word.
\follow ``Why did you choose it?'' I had a clear target and labelled data. Without labels, or to explore the data
first, I would choose unsupervised.
\end{qabox}

\begin{qabox}[Schlüter]{\asked Binary classification with labels 0 and 1: what is the output, why a sigmoid, and could we use a step function?}
\from{DL Basic Techniques (DLNN I) · Ch.~4 Loss functions, Ch.~6 Backprop · asked October 2026}
\ans The output is the probability of class 1. The sigmoid gives a smooth value in $(0,1)$, and with
cross-entropy its delta is prediction minus label. A step function has zero derivative almost everywhere, so
every delta is zero and nothing is learned.
\formula Step 1, the slope of the sigmoid:
\[ \sigma(s)=\frac{1}{1+e^{-s}},\qquad \sigma'(s)=\frac{e^{-s}}{(1+e^{-s})^2}=\sigma(s)\bigl(1-\sigma(s)\bigr)\le\tfrac14 . \]
Step 2, binary cross-entropy and its derivative:
\[ L=-\bigl[y\log a+(1-y)\log(1-a)\bigr],\qquad \frac{\partial L}{\partial a}=\frac{a-y}{a(1-a)} . \]
Step 3, the chain rule. The factors $a(1-a)$ cancel:
\[ \delta=\frac{\partial L}{\partial a}\,\sigma'(s)=a-y . \]
Step function: $H'(s)=0$ for $s\ne0$ and undefined at $0$, so $\delta=0$. The perceptron used a step, but with its
own rule, not backpropagation.
\pic{A flat staircase: standing on a step you cannot feel which way is down.}
\end{qabox}

\begin{qabox}[Schlüter]{\asked How does the learning rate influence training? Why small, why large?}
\from{DL Basic Techniques (DLNN I) · Ch.~7 Optimization · asked October 2026}
\ans The learning rate is the step size: too small and training crawls, too large and the loss oscillates or
diverges.
\formula Near a minimum the loss along one direction is a parabola with curvature $h$:
\[ L(w)=\tfrac12 h w^2,\quad \nabla L=hw,\quad w^{\text{new}}=w-\eta h w=(1-\eta h)\,w . \]
Each step multiplies the distance by $1-\eta h$. Convergence needs $|1-\eta h|<1$, so
\[ 0<\eta<\frac{2}{h}. \]
Tiny $\eta$: factor close to 1, slow. $\eta>2/h$: factor above 1 in size, divergence. In practice: warm-up, then
a decay (cosine).
\end{qabox}

\begin{qabox}[Schlüter]{\asked Learning rate versus batch size: how are they related?}
\from{DL Basic Techniques (DLNN I) · Ch.~7 · asked October 2026}
\ans A bigger batch gives a less noisy gradient, so it can take a bigger step: the linear scaling rule multiplies
the learning rate by the same factor as the batch size.
\formula Mini-batch gradient: unbiased, with variance divided by $B$:
\[ \mathbb{E}[\mathbf{g}_B]=\nabla\Remp,\qquad \mathrm{Var}[\mathbf{g}_B]=\frac{\sigma^2}{B}. \]
$k$ small steps with batch $B$ move about $\eta\sum_{t=1}^{k}\mathbf{g}_B^{(t)}\approx k\eta\,\mathbf{g}_{kB}$ (if the
gradient changes little), which is one step with batch $kB$ and rate $k\eta$. So
\[ B\to kB\quad\Longrightarrow\quad \eta\to k\eta \]
(Goyal et al.\ 2017, with a warm-up). The noise level is set by $\eta/B$. Very large batches stop paying off.
\end{qabox}

\begin{qabox}[Schlüter]{\asked Gradient descent versus Adam on mini-batches: how does each one use the data?}
\from{DL Basic Techniques (DLNN I) · Ch.~7 Adam · asked October 2026}
\ans Both take one step per mini-batch, but SGD uses only the current gradient with one global rate, while Adam
keeps running averages of past gradients and gives every weight its own step size.
\formula Adam, with $\mathbf{g}$ the gradient of the current batch:
\[ \mathbf{m}\leftarrow\beta_1\mathbf{m}+(1-\beta_1)\mathbf{g},\qquad \mathbf{v}\leftarrow\beta_2\mathbf{v}+(1-\beta_2)\mathbf{g}^2 . \]
Bias correction: both start at zero. Unrolling gives $\mathbb{E}[\mathbf{m}_t]=(1-\beta_1^t)\,\mathbb{E}[\mathbf{g}]$, so
\[ \hat{\mathbf{m}}=\frac{\mathbf{m}_t}{1-\beta_1^t},\quad \hat{\mathbf{v}}=\frac{\mathbf{v}_t}{1-\beta_2^t},\quad \mathbf{w}\leftarrow\mathbf{w}-\eta\,\frac{\hat{\mathbf{m}}}{\sqrt{\hat{\mathbf{v}}}+\epsilon}. \]
Scale invariance: $\mathbf{g}\to c\,\mathbf{g}$ gives $\hat{\mathbf{m}}\to c\,\hat{\mathbf{m}}$ and $\sqrt{\hat{\mathbf{v}}}\to c\sqrt{\hat{\mathbf{v}}}$, so the step is unchanged.
\pic{SGD forgets the last batch. Adam remembers.}
\end{qabox}

\begin{qabox}[Hochreiter]{\asked What is the difference between AdamW and Adam with L2 regularization?}
\from{DL Basic Techniques (DLNN I) · Ch.~7--8 · asked October 2026}
\ans AdamW applies the weight decay directly to the weights, outside the adaptive division. Adam with L2 sends the
decay through the division, so weights with large gradients are hardly decayed.
\formula For SGD the two are the same:
\[ \mathbf{w}-\eta(\mathbf{g}+\lambda\mathbf{w})=(1-\eta\lambda)\,\mathbf{w}-\eta\,\mathbf{g}. \]
Adam with L2 uses $\mathbf{g}+\lambda\mathbf{w}$ inside the moments, so the decay is divided by $\sqrt{\hat{\mathbf{v}}}$. AdamW
(Loshchilov and Hutter 2019) keeps it outside:
\[ \mathbf{w}\leftarrow\mathbf{w}-\eta\Bigl(\frac{\hat{\mathbf{m}}}{\sqrt{\hat{\mathbf{v}}}+\epsilon}+\lambda\mathbf{w}\Bigr). \]
\end{qabox}

\begin{qabox}[Hochreiter]{\asked Which normalizations do you know, and which one is popular now?}
\from{DL Basic Techniques (DLNN I) · Ch.~10 Normalization · asked October 2026}
\ans Batch norm normalizes each feature over the batch, layer norm normalizes over the features of one sample,
and RMSNorm, the popular one in large language models, only divides by the root mean square.
\formula One recipe for a set of numbers $x_1,\dots,x_d$:
\[ \hat x_i=\frac{x_i-\mu}{\sqrt{\sigma^2+\epsilon}},\qquad y_i=\gamma\hat x_i+\beta . \]
Batch norm: the set is one feature over the batch. Layer norm: all features of one sample. RMSNorm (Zhang and
Sennrich 2019) drops the mean:
\[ y_i=\frac{x_i}{\sqrt{\frac1d\sum_k x_k^2+\epsilon}}\,g_i . \]
The rescaling, not the re-centring, stabilizes training, and RMSNorm is cheaper (used in LLaMA).
\end{qabox}

\begin{qabox}[Hochreiter]{\asked Which initialization techniques do you know? Why this variance, and what is the theory?}
\from{DL Basic Techniques (DLNN I) · Ch.~10 Initialization · asked October 2026}
\ans LeCun, Glorot and He choose the weight variance so that the variance of the signal, and of the deltas, stays
constant from layer to layer.
\formula Independent, zero-mean weights, independent of the inputs:
\[ s=\sum_{j=1}^{n_{\text{in}}}w_ja_j\;\Longrightarrow\;\mathrm{Var}(s)=n_{\text{in}}\,\mathrm{Var}(w)\,\mathbb{E}[a^2]. \]
Constant variance needs $\mathrm{Var}(w)=1/n_{\text{in}}$ (LeCun, used with SELU). ReLU gives
$\mathbb{E}[a^2]=\tfrac12\mathrm{Var}(s)$, so $\mathrm{Var}(w)=2/n_{\text{in}}$ (He). The backward pass sums over
$n_{\text{out}}$, so Glorot takes the compromise
\[ \mathrm{Var}(w)=\frac{2}{n_{\text{in}}+n_{\text{out}}}. \]
Zero initialization fails: all units stay identical (symmetry). Orthogonal initialization sets all singular values to 1.
\end{qabox}

\begin{qabox}[Hochreiter]{\asked How is the spectral norm defined? Which structure do we want for $W$, and how should it change over learning?}
\from{DL Basic Techniques (DLNN I) · Ch.~6 Vanishing gradient; SS25 theory · asked October 2026}
\ans The spectral norm is the largest singular value, the largest factor by which $W$ can stretch a vector. We want
$W$ close to an isometry, with singular values near one, and the norm should stay bounded during training.
\formula Definition and its link to eigenvalues:
\[ \|W\|_2=\max_{\mathbf{x}\ne0}\frac{\|W\mathbf{x}\|}{\|\mathbf{x}\|},\qquad \|W\|_2^2=\max_{\mathbf{x}\ne0}\frac{\mathbf{x}^\top W^\top W\mathbf{x}}{\mathbf{x}^\top\mathbf{x}}=\lambda_{\max}(W^\top W)=\sigma_{\max}^2 . \]
With 1-Lipschitz activations, every layer stretches a difference at most by $\|W_l\|_2$, so
\[ \|f(\mathbf{x})-f(\mathbf{x}')\|\le\prod_l\|W_l\|_2\,\|\mathbf{x}-\mathbf{x}'\|. \]
If all $\sigma_i=1$, $W$ is orthogonal: signals and deltas keep their size. During training the weights learn
structure and their norms grow, so we keep them bounded (weight decay, spectral normalization $W/\|W\|_2$).
\end{qabox}

\begin{qabox}[Hochreiter]{\asked Explain vanishing gradients. What can we do to mitigate them?}
\from{DL Basic Techniques (DLNN I) · Ch.~6; LSTM \& Transformers (Hochreiter) · asked October 2026}
\ans The delta of an early layer is a product of one factor per layer, so it vanishes or explodes when these factors
are below or above one. Every fix keeps the factors close to one.
\formula Unroll the delta recursion and take norms:
\[ \boldsymbol\delta^{(l)}=f'(\mathbf{s}^{(l)})\odot W^{(l+1)\top}\boldsymbol\delta^{(l+1)},\qquad
\|\boldsymbol\delta^{(1)}\|\le\prod_{l}\Bigl(\max|f'|\;\|W^{(l)}\|_2\Bigr)\|\boldsymbol\delta^{(L)}\| . \]
Sigmoid: $\max|f'|=\tfrac14$, so if $\tfrac14\|W\|_2<1$ the delta decays exponentially with depth (Hochreiter 1991).
Residual block: $\mathbf{a}^{\text{new}}=\mathbf{a}+F(\mathbf{a})$ gives the Jacobian $I+\partial F/\partial\mathbf{a}$, which does not vanish.
Fixes: initialization, ReLU/ELU/SELU, normalization, residual connections, gradient clipping, LSTM.
\end{qabox}

\begin{qabox}[Hochreiter]{\asked Write down the SVM primal. What do we optimize, where does the ``1'' come from, and could we use another value? What is the canonical form?}
\from{ML: Supervised Techniques (Lehner) · SVM · asked October 2026}
\ans We maximise the margin. The 1 only fixes the free scale of $(\mathbf{w},b)$, so any positive value gives the same
hyperplane. Fixing it to 1 is the canonical form.
\formula Distance of a point to the hyperplane, and the free scale:
\[ d_i=\frac{|\mathbf{w}^\top\mathbf{x}_i+b|}{\|\mathbf{w}\|},\qquad (\mathbf{w},b)\to(c\mathbf{w},cb)\ \text{gives the same hyperplane.} \]
Canonical form: $\min_i y_i(\mathbf{w}^\top\mathbf{x}_i+b)=1$, so the margin is $2/\|\mathbf{w}\|$. Maximising it means
\[ \min_{\mathbf{w},b}\ \tfrac12\|\mathbf{w}\|^2\quad\text{s.t.}\quad y_i(\mathbf{w}^\top\mathbf{x}_i+b)\ge1\ \ \forall i . \]
With $c$ instead of 1, put $\mathbf{w}'=\mathbf{w}/c$, $b'=b/c$: the constraint becomes $\ge1$ again. Same classifier.
\end{qabox}

\begin{qabox}[Schlüter]{\asked CNN: what is the receptive field, how can we increase it, and what does pooling do?}
\from{DL Basic Techniques (DLNN I) · CNN chapter · asked October 2026}
\ans The receptive field is the part of the input that can influence one unit. It grows with depth, kernel size,
stride or pooling, and dilation.
\formula A layer with kernel $k$ adds $k-1$ positions, measured in the jump $j$ (the product of earlier strides):
\[ r_l=r_{l-1}+(k_l-1)\,j_{l-1},\qquad j_l=j_{l-1}s_l,\qquad r_0=j_0=1 . \]
Three $3\times3$ layers, stride 1: $r=1+2+2+2=7$. Stride or pooling 2 doubles $j$. Dilation $d$: effective kernel
$d(k-1)+1$. Pooling (max or average) has no weights, shrinks the map, enlarges later receptive fields and adds a
little shift invariance. Global average pooling gives one vector per image or sequence.
\end{qabox}
